% problem-000466 7 环氧氯化物开环选择性
\section*{7. 环氧氯化物开环选择性}
海洋⽣物体内蕴含着丰富的⽣物活性物质，误⻝某些海鲜⻝品会导致严重的⻝物中毒。细胞毒素chlorosulpholipid是从某海洋⽣物体中分离得到的⼀种天然产物，其全合成研究表明，该分⼦内含有的多氯代结构对其⽣物活性具有重要意义。在该分⼦合成路线中，氯原⼦的引⼊涉及如下所⽰的环氧开环反应 $\mathrm { ( T B S = } ~ t - \mathrm { ) }$ $\mathrm { B u ( C H _ { 3 } ) _ { 2 } S i ) }$ ：

（图：）

以⼆氯甲烷与⼄酸⼄酯为溶剂，得到3种开环产物 $\mathrm { \bf P } _ { 1 } , \mathrm { \bf P } _ { 2 }$ 和 ${ \bf P } _ { 3 }$ 。

\subsection*{7-1}
给出中间体M的结构简式。

\subsection*{7-2}
给出 $\mathbf { P } _ { 1 } , \ \mathbf { P } _ { 2 }$ 和 ${ \bf P } _ { 3 }$ 的结构简式，以及由M得到相应产物的⽣成机理。
